\documentclass[conference]{IEEEtran}
\usepackage{times}

\usepackage[numbers]{natbib}
\usepackage{multicol}
\usepackage{graphicx}
\usepackage{booktabs}
\usepackage{multirow}
\usepackage{xcolor}
\let\labelindent\relax
\usepackage{enumitem}

\usepackage[bookmarks=true]{hyperref}

\pdfinfo{
   /Author (Anonymous)
   /Title  (AutoDex Rebuttal)
   /CreationDate (D:20260327120000)
   /Subject (Robotics)
   /Keywords (Dexterous;Grasping)
}

\definecolor{R1}{RGB}{0,128,0}
\definecolor{R2}{RGB}{0,0,200}
\definecolor{R3}{RGB}{200,0,0}
\definecolor{AC}{RGB}{128,0,128}

\newcommand{\rone}{{\color{R1}R1}}
\newcommand{\rtwo}{{\color{R2}R2}}
\newcommand{\rthree}{{\color{R3}R3}}
\newcommand{\acmark}{{\color{AC}AC}}

\renewcommand{\thetable}{R-\arabic{table}}
\renewcommand{\thefigure}{R-\arabic{figure}}

\newcommand{\rssparagraph}[1]{{\bf #1.}}

% Allow up to 4 floats per column top, using up to 95% of column
\setcounter{topnumber}{4}
\renewcommand{\topfraction}{0.95}
\setcounter{totalnumber}{4}

\begin{document}

\title{AutoDex [Paper-ID: 99]\vspace{-1cm}}
\maketitle
\thispagestyle{empty}

%% ============================================================
%% TEXT BODY
%% ============================================================

We thank the reviewers for recognizing the importance of bridging the sim-to-real gap ({\color{R1}R1}, {\color{R3}R3}), the ``clear quantitative argument for real-world filtering'' and ``coverage-preserving selection design'' ({\color{R2}R2}), and the inclusion of both successful and failed grasps ({\color{R3}R3}). We address all concerns below.

%% --- SECTION 1: Geometric Transfer ---

\rssparagraph{[\rthree, \rtwo, \acmark] Geometric transfer and grasp selection}
Dexterous grasping in constrained environments requires collision-free motion planning \emph{and} stable object contact. The former --- which we call \emph{geometric feasibility} (IK solvability, collision avoidance) --- transfers reliably from simulation to the real world because robot kinematics and obstacle geometry are known precisely. We exploit this property: grasps that are geometrically feasible in simulation remain feasible on the real robot.

[\rthree] To systematically cover constrained scenarios, we decompose environmental constraints into planar boundaries: one side = Wall; one side + ceiling = Shelf; 2--3 sides $\pm$ ceiling = Shelf variants; 4 sides open-top = Box. The only excluded case --- fully enclosed --- is infeasible for any method. Given these constraints, we select grasps via a coverage-preserving greedy algorithm that maximizes solvable scenarios per real-robot trial.

[\rtwo] Table~\ref{tab:selection} confirms that greedy selection outperforms random at practical budgets ($N$=10--100), with the largest gain at small $N$ where each trial is most valuable. At large $N$, random sampling achieves sufficient coverage through volume, so convergence is expected. The sim solvability ceiling converges at $\sim$[X]\% around $N$=[X]; of the remaining unsolved scenarios, [X]\% are kinematically infeasible (obstacle density blocks all approach directions) and only [X]\% are coverage gaps.

%% --- SECTION 2: Contact doesn't transfer ---

\rssparagraph{[\rone, \rthree, \acmark] Contact dynamics do not transfer: the necessity of real-world validation}
Unlike geometry, contact dynamics are \emph{not} reliably simulated. DR randomizes known parameters (friction, mass), but the Allegro hand's dominant failure modes --- finger-pad deformation, cable backlash, non-linear contact stiffness --- are not parameterized in current simulators.

Table~\ref{tab:dr_fp_fn} quantifies this gap on 227 grasps: 83.8\% of real-world successful grasps were rejected by sim DR (false negatives --- valid grasps permanently discarded), while 7 of 18 sim-approved grasps (38.9\%) still failed on the real robot. This demonstrates that sim DR is simultaneously too conservative (discarding the majority of viable grasps) and insufficiently reliable (nearly 40\% of its approvals fail in practice).

[\rthree, \acmark] We acknowledge the conceptual tension: we claim contact doesn't transfer, yet use simulation filtering. The key is \emph{how} we use it. Our sim filter operates under deliberately optimistic settings (high friction, low gravity) to remove only obvious failures --- grasps with zero contact or mesh penetration --- while preserving all borderline candidates. This is a coarse pre-filter, not a contact-dynamics oracle. Real execution filters out the false positives this check misses.

%% --- SECTION 3: Physical filtering gains ---

\rssparagraph{[\rtwo, \rthree, \acmark] Performance gains through physical filtering}
We expanded evaluation with more objects, more trials, and confidence intervals (Fig.~\ref{fig:physical_validation}). Real-filtered grasps achieve [X]\% $\pm$ [X] vs [X]\% $\pm$ [X] for sim-validated (Fisher's exact $p < $ [X]), with consistent gains across all material and scene categories. This confirms that physical filtering captures dynamics that simulation cannot predict.

%% ============================================================
%% FLOAT — single figure* spanning both columns
%% ============================================================

% --- Table R-1: Both tables stacked ---
\begin{table}[t]
\caption{(a) Scenario solvability (\%) by grasp budget. (b) Sim DR vs real ($N$=227).}
\label{tab:selection}
\vspace{-4pt}
\scriptsize
\begin{minipage}[t]{0.48\columnwidth}
\centering
{\small\textbf{(a)}}
\vspace{1pt}

\begin{tabular}{l|ccc}
\toprule
& \!$N$\!=\!10\! & \!$N$\!=\!50\! & \!$N$\!=\!100\! \\
\midrule
Random & & & \\
Greedy & & & \\
\bottomrule
\end{tabular}
\end{minipage}%
\hfill
\begin{minipage}[t]{0.48\columnwidth}
\centering
{\small\textbf{(b)}}
\label{tab:dr_fp_fn}
\vspace{1pt}

\begin{tabular}{l|cc}
\toprule
& \textbf{R.Succ} & \textbf{R.Fail} \\
\midrule
S.Pass & 11\,(TP) & 7\,(FP) \\
S.Fail & 57\,(FN) & 152\,(TN) \\
\bottomrule
\end{tabular}

\vspace{1pt}
{\scriptsize FPR:4.4\% FNR:\textbf{83.8\%}}
\end{minipage}
\end{table}

% --- Figure R-1: Both figures stacked ---
\begin{figure}[t]
\centering
\includegraphics[width=\columnwidth]{figures/physical_validation.pdf}
\vspace{-6pt}

{\small\textbf{(a)} Physical validation + object set}
\label{fig:physical_validation}

\vspace{2pt}
\includegraphics[width=\columnwidth]{figures/cost_combined.png}
\vspace{-6pt}

{\small\textbf{(b)} Collection efficiency \& camera trade-off}
\label{fig:cost_throughput}
\vspace{-4pt}
\caption{(a) Physical validation across object categories. (b) Cost and camera trade-offs.}
\end{figure}

%% ============================================================
%% TEXT BODY (continued)
%% ============================================================

%% --- SECTION 4: Teleop comparison ---

\rssparagraph{[\rtwo, \acmark] Efficiency compared to teleoperation}
Two methods exist for real-world grasp validation: teleoperation and our automated pipeline. AutoDex achieves [X] grasps/hr vs [X] for teleop, requiring only [X] min of human time per grasp (object placement). Teleop requires continuous operator presence (body suit, 22-DoF control), with sessions limited to [X] hours before fatigue and Xsens charging (Fig.~\ref{fig:cost_throughput}). AutoDex runs unattended overnight and scales linearly with robots without additional operators.

Beyond efficiency, AutoDex provides \emph{diversity}: teleop datasets are biased toward operator-preferred approach paths, and constructing constrained scenes (wall, shelf) during teleop is impractical. AutoDex evaluates grasps under procedural constraints via geometric transfer, and is generation-method agnostic --- any generator (BODex, RL, human priors) can be plugged in. We focus on stable grasps in this work --- an intentional scope decision, as no physically-validated dexterous grasp dataset existed even for this foundational task. The pipeline extends to functional grasps by swapping the generation method; validation and selection remain unchanged.

%% --- SECTION 5: Camera scalability ---

\rssparagraph{[\rthree, \acmark] Scalability of the multi-camera setup}
The 24-camera setup is a research choice, not a pipeline requirement. Our system is agnostic to camera count and type. Calibration: $\sim$40 board images ($\sim$10 min) + automated hand-eye ($\sim$10 min), regardless of camera count. During grasping, the hand occludes [X]\% of the object on average. More cameras mitigate this, but the trade-off is per-lab (Table~III: 4-cam ADD-S 1.17mm / 51\% success vs 24-cam 0.31mm / 62\%).

{
  \small
  \bibliographystyle{plainnat}
  \bibliography{references}
}

\end{document}
